\documentclass[12pt,a4paper]{article}
\usepackage[T2A]{fontenc}
\usepackage[utf8]{inputenc}
\usepackage[russian]{babel}
\usepackage{amsmath,amssymb,booktabs}
\usepackage[margin=25mm]{geometry}
\usepackage{hyperref}
\title{Задание 12\\Прямоугольные сечения со стороной $b(h)=b_0h^k$}
\author{}
\date{}
\begin{document}
\maketitle

\section{Условие}
Горизонтальные сечения сосуда являются прямоугольниками. Одна сторона имеет постоянную длину $a$, другая зависит от высоты: $b(h)=b_0h^k$. Требуется построить и исследовать модель вытекания воды через небольшое отверстие у дна.
Обозначим площадь отверстия через $\sigma$, ускорение свободного падения через $g$, начальную высоту воды через $H$. Предполагаем $a,b_0,\sigma,g,H>0$, $k\in\mathbb R$, $h(0)=H$. Высота отсчитывается от отверстия. Используется идеальная формула Торричелли без поправочного коэффициента расхода:
\[
v(h)=\sqrt{2gh}.
\]
Давление над водой и на выходе одинаково; потерями энергии пренебрегаем. Закон рассматривается при $h>0$. Показатель $k$ относится к одной стороне прямоугольника, поэтому площадь пропорциональна $h^k$, а не $h^{2k}$.

\section{Построение дифференциального уравнения}
Площадь сечения равна
\[
S(h)=ab(h)=ab_0h^k.
\]
Баланс объёма на малом промежутке времени даёт
\[
S(h)\,dh=-\sigma\sqrt{2gh}\,dt.
\]
Следовательно,
\begin{equation}
\boxed{ab_0h^k\frac{dh}{dt}=-\sigma\sqrt{2g}\,h^{1/2},\qquad h(0)=H.}
\end{equation}
Введём обозначения
\[
A=\frac{\sigma\sqrt{2g}}{ab_0}>0,\qquad p=k+\frac12.
\]
Тогда задача Коши принимает вид
\begin{equation}
\boxed{h'=-Ah^{1/2-k}=-Ah^{1-p},\qquad h(0)=H.}
\end{equation}
При $h>0$ правая часть гладкая, поэтому положительное решение единственно. Кроме того, $h'<0$: уровень строго убывает.

\section{Неявное и явное решения}
\subsection{Случай $k\ne-\tfrac12$}
Разделяя переменные, получаем
\[
h^{k-1/2}\,dh=-A\,dt,\qquad
\frac{h^{k+1/2}}{k+1/2}=C-At.
\]
Это общее решение в неявной форме. Из условия $h(0)=H$ следует $C=H^p/p$. Таким образом, неявное решение начальной задачи:
\begin{equation}
\boxed{\frac{h^p-H^p}{p}+At=0,\qquad p\ne0.}
\end{equation}
В исходных обозначениях:
\[
\frac{ab_0}{k+1/2}\bigl(h^{k+1/2}-H^{k+1/2}\bigr)
=-\sigma\sqrt{2g}\,t.
\]
Явное решение и обратная зависимость имеют вид
\begin{equation}
\boxed{h(t)=\bigl(H^p-pAt\bigr)^{1/p}},\qquad
\boxed{t(h)=\frac{H^p-h^p}{pA}}.
\end{equation}
Выбирается положительная вещественная ветвь: $H^p-pAt>0$. Если $p>0$, это условие выполняется до опорожнения; если $p<0$, оно выполняется при всех $t\ge0$.

\subsection{Случай $k=-\tfrac12$}
Теперь $h'=-Ah$. Интегрирование даёт $\ln h=C-At$, где константа зависит от выбранной единицы высоты. Удобная размерностно корректная форма начального решения:
\begin{equation}
\boxed{\ln\frac{h}{H}+At=0},\qquad
\boxed{h(t)=H e^{-At}},\qquad
\boxed{t(h)=\frac1A\ln\frac{H}{h}}.
\end{equation}
Уровень стремится к нулю, но не достигает нуля за конечное время.

\section{Время полного вытекания}
Время опорожнения определяется несобственным интегралом
\[
T=\frac1A\int_0^H h^{k-1/2}\,dh.
\]
Он сходится тогда и только тогда, когда $k-1/2>-1$. Поэтому
\begin{equation}
\boxed{T(H)=\begin{cases}
\displaystyle\frac{ab_0H^{k+1/2}}{\sigma\sqrt{2g}(k+1/2)},&k>-1/2,\\[6pt]
+\infty,&k\le-1/2.
\end{cases}}
\end{equation}
При конечном $T$ решение можно записать как
\[
h(t)=H\left(1-\frac{t}{T}\right)^{1/(k+1/2)},\qquad 0\le t<T.
\]
Значение $h(T)=0$ понимается как предел. После опорожнения физически полагают $h(t)=0$, не продолжая степенную формулу за $T$. Это отдельное описание пустого сосуда, а не утверждение о классической разрешимости исходного уравнения в точке $h=0$.
Для фиксированных $k>-1/2$ и остальных параметров
\[
\frac{dT}{dH}=\frac{H^{k-1/2}}A>0.
\]
Время увеличивается с начальной высотой, пропорционально $a,b_0$ и обратно пропорционально $\sigma\sqrt g$.

\section{Исследование интегральных кривых}
Дифференцируя уравнение при $h>0$, находим
\[
h''=A^2\left(\frac12-k\right)h^{-2k}.
\]
\begin{center}
\begin{tabular}{lll}
\toprule
Показатель & Изменение скорости уровня & Кривизна\\
\midrule
$k<1/2$ & $|h'|\to0$ при $h\to0$ & $h''>0$\\
$k=1/2$ & $h'=-A$ постоянно & $h''=0$\\
$k>1/2$ & $|h'|\to\infty$ при $h\to0$ & $h''<0$\\
\bottomrule
\end{tabular}
\end{center}
Следовательно, равномерная шкала водяных часов получается при
\[
\boxed{k=\frac12,\qquad b(h)=b_0\sqrt h,\qquad h(t)=H-At.}
\]
Необходимо различать скорость истечения $v=\sqrt{2gh}$ и скорость снижения уровня $|h'|$: первая стремится к нулю во всех случаях, а вторая зависит также от площади сечения.

\subsection{Основные частные случаи}
Для краткости используем ранее введённый коэффициент $A$; при изменении $k$ размерность $b_0$ и $A$ меняется.
\begin{center}
\begin{tabular}{lll}
\toprule
$k$ & $h(t)$ до опорожнения & $T$\\
\midrule
$0$ & $(\sqrt H-At/2)^2$ & $2\sqrt H/A$\\
$1/2$ & $H-At$ & $H/A$\\
$1$ & $(H^{3/2}-3At/2)^{2/3}$ & $2H^{3/2}/(3A)$\\
$2$ & $(H^{5/2}-5At/2)^{2/5}$ & $2H^{5/2}/(5A)$\\
\bottomrule
\end{tabular}
\end{center}
В отличие от конуса вращения, при $k=1$ здесь $S\sim h$, а не $h^2$.

\subsection{Отрицательные показатели и объём}
При $k<0$ ширина $b(h)$ неограниченно увеличивается при приближении к дну. Поэтому такая форма является математической идеализацией.
Объём жидкости конечен только при $k>-1$:
\[
V(H)=\int_0^H ab_0z^k\,dz=\frac{ab_0}{k+1}H^{k+1},\qquad k>-1.
\]
Имеются следующие режимы:
\begin{itemize}
\item $-1/2<k<0$: ширина у дна неограниченна, но объём и время вытекания конечны.
\item $k=-1/2$: конечный объём, экспоненциальное убывание уровня, $T=\infty$.
\item $-1<k<-1/2$: конечный объём, но бесконечное время вытекания. Если $q=-k-1/2>0$, то
\[
h(t)=\bigl(H^{-q}+qAt\bigr)^{-1/q}\sim(qAt)^{-1/q}.
\]
\item $k\le-1$: объём от нуля до любой положительной высоты бесконечен. Формула для положительного уровня остаётся формальным решением локального баланса, но не описывает обычный сосуд с конечным запасом воды.
\end{itemize}

\section{Проверка решения}
Для $p\ne0$ положим $B(t)=H^p-pAt>0$. Тогда
\[
h'=\frac1pB^{1/p-1}(-pA)=-AB^{(1-p)/p}=-Ah^{1-p},\qquad h(0)=H.
\]
Для $p=0$ непосредственно $(He^{-At})'=-AHe^{-At}$. Обе формулы удовлетворяют начальной задаче на области $h>0$.

\section{Программная реализация}
Программа состоит из двух файлов: \texttt{main.py} задаёт параметры и запускает расчёт, \texttt{model.py} содержит формулы, численные проверки, табличный вывод и построение графиков. Используются NumPy и Matplotlib. Решение вычисляется аналитически.
В демонстрации выбраны, а не взяты из условия, параметры $a=1$ м, $H=1$ м, ширина верхнего сечения $B=1$ м, $\sigma=10^{-3}$ м$^2$, $g=9.8$ м/с$^2$. Для сравнения форм задаётся
\[
b_0=\frac{B}{H^k},\qquad b(h)=B(h/H)^k.
\]
Так все сосуды имеют одинаковую верхнюю ширину. Они не имеют одинакового объёма. Параметр $b_0$ в исходном законе имеет размерность $\text{м}^{1-k}$, поэтому нельзя сравнивать разные $k$, считая его одной и той же размерной длиной.
При таком сравнении
\[
T=\frac{aB\sqrt H}{\sigma\sqrt{2g}(k+1/2)},\quad k>-1/2.
\]
Если при изменении $H$ сохранять верхнюю ширину $B$, меняется сам сосуд (коэффициент $b_0$). Это не та же зависимость, что $T(H)$ для фиксированного $b_0$.
Программа выводит времена и объёмы, проверяет $h(0)=H$ и согласованность $t(h(t))=t$ до опорожнения. Показываются графики $h(t)$ для разных $k$, семейство решений при разных начальных высотах в одном фиксированном сосуде, обратная зависимость $t(h)$ и зависимость $T(H)$ для фиксированного $b_0$. Графики показываются в окне; сохранять файлы для работы программы не требуется.

\section{Границы применимости и вывод}
Модель Торричелли предполагает малое отверстие по сравнению с сечением сосуда и пренебрежимо малую скорость свободной поверхности. При $k>0$ площадь $S(h)$ стремится к нулю, поэтому для фиксированного отверстия эти предположения нарушаются достаточно близко к дну. Неограниченный рост $|h'|$ при $k>1/2$ является признаком ограничения модели, а не физически достижимой бесконечной скорости.

Итак, явное и неявное решения найдены для всех вещественных $k$ при положительном уровне. Конечное время опорожнения получается при $k>-1/2$, постоянная скорость снижения уровня --- при $k=1/2$, конечный начальный объём --- при $k>-1$. Эти три условия отвечают на разные вопросы и не должны смешиваться.
\end{document}
